A first relevant literature studies how election administration and voting technology shape both political inclusion and the perceived legitimacy of elections. In Brazil, \citet{Fujiwara2015VotingTechnoloPolitica} shows that the adoption of electronic voting reduced residual votes, expanded de facto enfranchisement among poorer and less educated citizens, and increased political responsiveness, while \citet{Zucco2016TradingOldErrors} shows that the same technological transition also generated new forms of voting error. Those papers establish that administrative technology inside the voting process can materially affect representation, but they focus on ballot casting rather than voter identification. More broadly, work on election administration emphasizes that these procedures are politically consequential rather than a neutral background to democratic competition \citep{Ferrer2023HowPartisanLocal, Bubeck2017ProcessCandidatThe}. Relative to this literature, the present paper studies a later stage of electoral modernization and asks whether changing the registration and verification regime affected both who remained in the electorate and how citizens evaluated democratic institutions.

The paper also connects to the literature on participation costs and administrative burden. Classic work shows that registration laws materially affect turnout \citep{Rosenstone1978TheEffectRegistra, Ansolabehere2005TheIntroducVoter}, while later studies find that lowering pre-election frictions through preregistration or automatic reregistration can raise participation \citep{Holbein2014MakingYoungVoters, Kim2022AutomatiVoterReregist}. Related research on voter-identification rules shows that documentary requirements can shift participation across groups even when average turnout effects appear modest \citep{Fraga2021WhoVoterLaws}. This is central for the present project because biometric registration was designed as a security reform, but it also required pre-election compliance and therefore could change not only the size of the electorate but also its composition. In that respect, the paper also speaks to work showing that unequal turnout matters for political representation \citep{Hajnal2005WhereTurnoutMatters}.

Another adjacent literature studies biometrics, state capacity, and trust in public institutions. Outside elections, \citet{Muralidharan2016BuildingStateCapacity} and \citet{Bossuroy2024BiometriMonitoriService} show that biometric systems can improve targeting, monitoring, and administrative data quality by changing how citizens interact with the state. On the attitudinal side, \citet{Stevenson2011TrustPublicInstitut} shows that confidence in public institutions moves with broader economic conditions, while \citet{Marx2021VoterMobilisaAnd} provides direct evidence that mobilization can increase trust in electoral institutions. These papers make it plausible that a salient reform in election administration could have downstream effects on trust, but they also underscore that attitudinal responses are likely to be context dependent. For that reason, the trust and backlash evidence in this paper is best understood as a complementary extension to the administrative results rather than as the paper's core claim.

The closest evidence for Brazil is \citet{forquesatoRodrigues2023}, who study the turnout and registry effects of biometric recadastramento, and \citet{Karim2025VoterBuyingPolitici}, who examines a related re-registration reform aimed at curbing voter-buying and traces its consequences for political competition and public-good provision. Relative to that work, this paper combines municipality-level administrative outcomes with survey evidence on democratic attitudes, allowing it to study registry size, electorate composition, trust, and integrity beliefs within a common framework. Because rollout was staggered across municipalities, the paper also joins the recent literature that uses estimators designed for heterogeneous treatment timing rather than relying mechanically on conventional two-way fixed effects \citep{Callaway2020DifferenDifferenWith, Borusyak2024RevisitiEventStudy, deChaisemartin2020TwoWayFixed, GoodmanBacon2018DifferenDifferenWith}. Taken together, these contributions position the paper at the intersection of electoral modernization, administrative burden, and the political economy of trust.
